% Options for packages loaded elsewhere
\PassOptionsToPackage{unicode}{hyperref}
\PassOptionsToPackage{hyphens}{url}
\PassOptionsToPackage{dvipsnames,svgnames,x11names}{xcolor}
\documentclass[
]{article}
\usepackage{xcolor}
\usepackage[a4paper, hmargin=2.2cm, vmargin=2.4cm, headheight=14.5pt, footskip=0.9cm, includehead, includefoot]{geometry}
\usepackage{amsmath,amssymb}
\setcounter{secnumdepth}{-\maxdimen} % remove section numbering
\usepackage{iftex}
\ifPDFTeX
  \usepackage[T1]{fontenc}
  \usepackage[utf8]{inputenc}
  \usepackage{textcomp} % provide euro and other symbols
\else % if luatex or xetex
  \usepackage{unicode-math} % this also loads fontspec
  \defaultfontfeatures{Scale=MatchLowercase}
  \defaultfontfeatures[\rmfamily]{Ligatures=TeX,Scale=1}
\fi
\usepackage{lmodern}
\ifPDFTeX\else
  % xetex/luatex font selection
\fi
% Use upquote if available, for straight quotes in verbatim environments
\IfFileExists{upquote.sty}{\usepackage{upquote}}{}
\IfFileExists{microtype.sty}{% use microtype if available
  \usepackage[]{microtype}
  \UseMicrotypeSet[protrusion]{basicmath} % disable protrusion for tt fonts
}{}
\makeatletter
\@ifundefined{KOMAClassName}{% if non-KOMA class
  \IfFileExists{parskip.sty}{%
    \usepackage{parskip}
  }{% else
    \setlength{\parindent}{0pt}
    \setlength{\parskip}{6pt plus 2pt minus 1pt}}
}{% if KOMA class
  \KOMAoptions{parskip=half}}
\makeatother
\usepackage{longtable,booktabs,array}
\newcounter{none} % for unnumbered tables
\usepackage{calc} % for calculating minipage widths
% Correct order of tables after \paragraph or \subparagraph
\usepackage{etoolbox}
\makeatletter
\patchcmd\longtable{\par}{\if@noskipsec\mbox{}\fi\par}{}{}
\makeatother
% Allow footnotes in longtable head/foot
\IfFileExists{footnotehyper.sty}{\usepackage{footnotehyper}}{\usepackage{footnote}}
\makesavenoteenv{longtable}
\setlength{\emergencystretch}{3em} % prevent overfull lines
\providecommand{\tightlist}{%
  \setlength{\itemsep}{0pt}\setlength{\parskip}{0pt}}
% ============================================================
% 土力学速查手册 —— 学术黑白 A4 版式（xelatex / pandoc）
% 主题：严谨、黑白、无彩色、无装饰
% ============================================================
\usepackage[fontset=windows]{ctex}            % 中文排版（宋体正文/黑体标题，自动伪粗体）
% 把常用符号划入 CJK 字体（SimSun），避免落到拉丁字体缺字：
% 带圈数字 ①-⑩、几何符号 ◆■□、箭头 →、希腊字母 αφ、数学符号 ≈
\xeCJKDeclareCharClass{CJK}{"2460->"24FF}
\xeCJKDeclareCharClass{CJK}{"25A0->"25FF}
\xeCJKDeclareCharClass{CJK}{"2190->"21FF}
\xeCJKDeclareCharClass{CJK}{"0370->"03FF}
\xeCJKDeclareCharClass{CJK}{"2248->"2248}% ≈（SimSun 有；− 等保留拉丁字体）
\usepackage{amsmath, amssymb}                 % 数学符号
\usepackage{booktabs}                         % 三线表
\usepackage{longtable}                        % 跨页表格
\usepackage{array}
\usepackage{fancyhdr}
\usepackage{enumitem}
\setlist{itemsep=2pt, topsep=4pt, leftmargin=1.6em}

% ---------- 页眉页脚：简洁学术 ----------
\pagestyle{fancy}
\fancyhf{}
\fancyhead[C]{\small\kaishu 土力学速查手册}
\fancyfoot[C]{\small\thepage}
\renewcommand{\headrulewidth}{0.4pt}
\renewcommand{\footrulewidth}{0pt}
\fancypagestyle{plain}{
  \fancyhf{}
  \fancyfoot[C]{\small\thepage}
  \renewcommand{\headrulewidth}{0pt}
}

% ---------- 表格：黑白三线风格 ----------
\renewcommand{\arraystretch}{1.35}
\setlength{\aboverulesep}{2pt}
\setlength{\belowrulesep}{2pt}

% ---------- 其他 ----------
\raggedbottom
\usepackage{bookmark}
\IfFileExists{xurl.sty}{\usepackage{xurl}}{} % add URL line breaks if available
\urlstyle{same}
% fallback for those not using the hyperref driver hyperxmp:
\makeatletter
\@ifundefined{xmpquote}{\newcommand{\xmpquote}[1]{#1}}{}
\makeatother
\hypersetup{
  colorlinks=true,
  linkcolor={black},
  filecolor={Maroon},
  citecolor={black},
  urlcolor={black},
  pdfcreator={LaTeX via pandoc}}

\author{}
\date{}

\begin{document}

\section{第七章
土坡稳定分析}\label{ux7b2cux4e03ux7ae0-ux571fux5761ux7a33ux5b9aux5206ux6790}

\begin{quote}
\textbf{本章导览}：考纲集中于两类------\textbf{无黏性土坡}（简单公式
\(F_s=\tan\varphi/\tan\beta\)，必考）、\textbf{黏性土坡条分法}（瑞典条分法为主，简化毕肖普了解步骤）。期末常考「判断稳定否」与「条分法表格计算」；考研则以推导受力图、公式含义辨析为主。核心思想：\textbf{抗滑力（矩）与滑动力（矩）之比}。
\end{quote}

\begin{center}\rule{0.5\linewidth}{0.5pt}\end{center}

\subsection{7.1 基本概念}\label{ux57faux672cux6982ux5ff5}

\subsubsection{公式
7-1　安全系数定义（通用）}\label{ux516cux5f0f-7-1-ux5b89ux5168ux7cfbux6570ux5b9aux4e49ux901aux7528}

\[F_s=\frac{\text{抗滑力(矩)}}{\text{滑动力(矩)}}\]

\[F_s=\frac{\tau_{f}}{\tau}=\frac{R\text{（抗滑力矩）}}{M\text{（滑动力矩）}},\qquad F_s\ge1.0\text{ 稳定（工程常要求 }F_s\ge1.2\sim1.5\text{）}\]

\begin{center}\rule{0.5\linewidth}{0.5pt}\end{center}

\subsection{7.2
无黏性土坡（平面滑动）}\label{ux65e0ux9ecfux6027ux571fux5761ux5e73ux9762ux6ed1ux52a8}

\subsubsection{公式
7-2　无黏性土坡稳定系数}\label{ux516cux5f0f-7-2-ux65e0ux9ecfux6027ux571fux5761ux7a33ux5b9aux7cfbux6570}

\[\boxed{F_s=\frac{\tan\varphi}{\tan\beta}}\]

\begin{itemize}
\tightlist
\item
  \textbf{参数说明}：\(\varphi\) 为砂土内摩擦角；\(\beta\)
  为坡角；两者皆取（°）。
\item
  \textbf{推导}：沿坡面取土单元，下滑力 \(W\sin\beta\)，抗滑力
  \(W\cos\beta\tan\varphi\)，令其比即得------与坡高\textbf{无关}。
\item
  \textbf{判据}：\(\beta\le\varphi\)
  恒稳定（无论多高）；\(\beta>\varphi\) 破坏。暴雨/水位骤降时
  \(\varphi\) 下降（饱和后强度降低）导致滑塌------概念简答题素材。
\item
  \textbf{注意}：无黏性土坡稳定\textbf{只取决于坡角}，与坡高无关------这是与黏性土坡的本质区别，考试判断题高频。
\end{itemize}

\begin{center}\rule{0.5\linewidth}{0.5pt}\end{center}

\subsection{\texorpdfstring{7.3 黏性土坡整体圆弧滑动法（\(\varphi=0\)
法）}{7.3 黏性土坡整体圆弧滑动法（\textbackslash varphi=0 法）}}\label{ux9ecfux6027ux571fux5761ux6574ux4f53ux5706ux5f27ux6ed1ux52a8ux6cd5varphi0-ux6cd5}

\subsubsection{公式
7-3　整体圆弧滑动安全系数（力矩平衡）}\label{ux516cux5f0f-7-3-ux6574ux4f53ux5706ux5f27ux6ed1ux52a8ux5b89ux5168ux7cfbux6570ux529bux77e9ux5e73ux8861}

\[F_s=\frac{M_R}{M_S}\]

\textbf{更规范表达}：

\[F_s=\frac{c_u\,L\,R}{W\,d}=\frac{c_u}{W}\cdot\frac{L R}{d}\]

\begin{itemize}
\tightlist
\item
  \textbf{参数说明}：\(c_u\) 为不排水抗剪强度（黏性土饱和快剪
  \(\varphi_u=0\)）；\(L\) 为滑弧长度；\(R\) 为滑弧半径；\(W\)
  为滑动土体总重；\(d\) 为 \(W\) 作用线到圆心 O 的水平距离。
\item
  \textbf{注意}：本方法适用于均质黏性土、\(\varphi\approx0\)（饱和软黏土）的情形；当
  \(\varphi>0\) 时需条分法（见 7.4）。
\end{itemize}

\subsubsection{公式
7-4　稳定数（泰勒稳定数）}\label{ux516cux5f0f-7-4-ux7a33ux5b9aux6570ux6cf0ux52d2ux7a33ux5b9aux6570}

\[N_s=\frac{c}{\gamma H_{cr}}\]

\[H_{cr}=\frac{c}{\gamma N_s}=\frac{c\,N_{cr}}{\gamma}\quad(N_s\ \text{由坡角与}\ \varphi\ \text{查图})\]

\begin{itemize}
\tightlist
\item
  \textbf{参数说明}：\(H_{cr}\) 为临界（极限）坡高；\(N_s\)
  为稳定数（泰勒解曲线，按坡角与 \(\varphi\) 查）；\(c,\gamma\)
  为土强度与重度。
\item
  \textbf{用途}：已知坡形求极限坡高、或已知高度判断稳定（\(H\le H_{cr}\)）。
\item
  \textbf{注意}：\(N_s\) 随坡角增大而增大；对 \(\varphi=0\) 饱和黏土坡
  \(N_s\approx0.181\times\cdots\) 具体查表，考试以查图/给数为主。
\end{itemize}

\begin{center}\rule{0.5\linewidth}{0.5pt}\end{center}

\subsection{7.4
条分法（圆弧滑动法的一般情形）}\label{ux6761ux5206ux6cd5ux5706ux5f27ux6ed1ux52a8ux6cd5ux7684ux4e00ux822cux60c5ux5f62}

将滑动土体竖向分为若干土条（每条约宽 \(b_i\)），逐条受力分析后求和。

\subsubsection{公式
7-5　瑞典条分法（费兰纽斯法，忽略条间作用力）}\label{ux516cux5f0f-7-5-ux745eux5178ux6761ux5206ux6cd5ux8d39ux5170ux7ebdux65afux6cd5ux5ffdux7565ux6761ux95f4ux4f5cux7528ux529b}

\[\boxed{F_s=\frac{\sum\limits_{i=1}^{n}\Big[c_i\,l_i+\big(W_i\cos\alpha_i-u_i\,l_i\big)\tan\varphi_i\Big]}{\sum\limits_{i=1}^{n}W_i\sin\alpha_i}}\]

\begin{itemize}
\tightlist
\item
  \textbf{参数说明}：\(W_i\) 为第 \(i\) 条土条自重；\(\alpha_i\)
  为该条滑面中点切线与水平面夹角；\(l_i\) 为该条滑弧长度；\(u_i\)
  为滑面上的孔隙水压力；\(c_i,\varphi_i\) 为抗剪强度指标。
\item
  \textbf{分析}：分子为抗滑力总和（黏聚力项 +
  摩擦项），分母为下滑力总和；自重用 \(\sum W\sin\alpha\) 投影。
\item
  \textbf{注意}：\(\alpha_i\) 取度，切线与水平夹角------靠近坡顶的条
  \(\alpha\) 为正、坡底处为负；\(W\cos\alpha_i-u_il_i\)
  是扣除孔压后的有效法向力。
\item
  忽略条间力使结果\textbf{偏保守}（\(F_s\) 偏小）。
\end{itemize}

\subsubsection{公式
7-6　简化毕肖普法（条间力只考虑水平分量）}\label{ux516cux5f0f-7-6-ux7b80ux5316ux6bd5ux8096ux666eux6cd5ux6761ux95f4ux529bux53eaux8003ux8651ux6c34ux5e73ux5206ux91cf}

\[F_s=\frac{\sum\limits_{i=1}^{n}\dfrac{\Big[c_i' b_i+\big(W_i-u_i b_i\big)\tan\varphi_i'\Big]}{\cos\alpha_i\left(1+\dfrac{\tan\alpha_i\,\tan\varphi_i'}{F_s}\right)}}{\sum\limits_{i=1}^{n}W_i\sin\alpha_i}\]

\begin{itemize}
\tightlist
\item
  \textbf{参数说明}：\(b_i\)
  为条宽（\(b_i=l_i\cos\alpha_i\)）；\(c',\varphi'\)
  为\textbf{有效应力}指标；其余同前。
\item
  \textbf{特点}：上式右端含 \(F_s\)（隐式），需\textbf{迭代}：先设
  \(F_s\) 初值代入右端 → 得新值 → 反复至收敛（一般 2\textasciitilde3
  次）。
\item
  \textbf{注意}：简化毕肖普
  得到的结果比瑞典法略大（更接近严格解）；某些教材强调「简化毕肖普不满足所有静力平衡」但仍可接受。
\end{itemize}

\begin{center}\rule{0.5\linewidth}{0.5pt}\end{center}

\subsection{7.5
最危险滑动面与判别（图解法/经验法）}\label{ux6700ux5371ux9669ux6ed1ux52a8ux9762ux4e0eux5224ux522bux56feux89e3ux6cd5ux7ecfux9a8cux6cd5}

\begin{itemize}
\tightlist
\item
  \textbf{费尔纽斯最危险滑面法}：对于均质黏性土坡，最危险滑弧通过\textbf{坡脚}；圆心位于由坡角
  \(\beta\) 决定的\textbf{两条等角线交点}附近（\(\varphi=0\)
  时直接由两角定圆心，\(\varphi>0\) 时在交点的径向延长线上试算）。
\item
  \textbf{简化处理}：先定圆心范围 → 试算若干圆弧求 \(F_{s,min}\) →
  该面为最危险滑面。
\item
  \textbf{考试呈现}：只要求按给定滑面验算
  \(F_s\)，最危险面搜索属方法了解。
\end{itemize}

\begin{center}\rule{0.5\linewidth}{0.5pt}\end{center}

\subsection{7.6
渗流作用下的坡稳定性（概念）}\label{ux6e17ux6d41ux4f5cux7528ux4e0bux7684ux5761ux7a33ux5b9aux6027ux6982ux5ff5}

\subsubsection{公式
7-7　无黏性土坡顺坡渗流（渗流平行坡面）的安全系数}\label{ux516cux5f0f-7-7-ux65e0ux9ecfux6027ux571fux5761ux987aux5761ux6e17ux6d41ux6e17ux6d41ux5e73ux884cux5761ux9762ux7684ux5b89ux5168ux7cfbux6570}

\[F_s=\frac{\gamma'}{\gamma_{sat}}\cdot\frac{\tan\varphi}{\tan\beta}\]

\begin{itemize}
\tightlist
\item
  \textbf{参数说明}：\(\gamma',\gamma_{sat}\)
  为无黏性土的浮重度与饱和重度；\(\varphi,\beta\) 为内摩擦角与坡角。
\item
  \textbf{物理意义}：顺坡渗流时流线与坡面平行，下滑力按 \(\gamma_{sat}\)
  计、抗滑的有效法向力按 \(\gamma'\) 计，故比无渗流的
  \(F_s=\tan\varphi/\tan\beta\)（公式 7-2）多乘
  \(\gamma'/\gamma_{sat}\)（常为 0.5 左右）。判据
  \(\beta\le\arctan\left(\dfrac{\gamma'}{\gamma_{sat}}\tan\varphi\right)\)。
\item
  \textbf{其他渗流方向}：渗流方向改变时分母项随之变化------如渗流竖直向下时回到
  \(F_s=\tan\varphi/\tan\beta\)
  附近，水平向外渗流时一般比顺坡情形更不利；考试以\textbf{平行坡面（顺坡出流）}为最常见情形。
\item
  \textbf{注意}：此式是「无黏性土坡 +
  渗流」专题的核心式，务必能由「有效应力法受力分析」推导出来。
\end{itemize}

\subsubsection{渗流对条分法的影响}\label{ux6e17ux6d41ux5bf9ux6761ux5206ux6cd5ux7684ux5f71ux54cd}

\begin{itemize}
\tightlist
\item
  渗流（尤其\textbf{平行坡面向外流}）使坡内孔隙水压力升高、有效应力降低，\(F_s\)
  下降。
\item
  处理：条分法中 \(u_i\)
  计入孔压（有效应力法）；极端水压（暴雨水位骤升/骤降）是滑坡诱因。
\item
  工程措施：排水沟/盲沟、削坡减载、坡脚反压、挡土结构------「措施题」备三类。
\end{itemize}

\begin{center}\rule{0.5\linewidth}{0.5pt}\end{center}

\subsection{7.7
本章公式索引小结}\label{ux672cux7ae0ux516cux5f0fux7d22ux5f15ux5c0fux7ed3}

{\def\LTcaptype{none} % do not increment counter
\begin{longtable}[]{@{}
  >{\raggedright\arraybackslash}p{(\linewidth - 4\tabcolsep) * \real{0.3333}}
  >{\raggedright\arraybackslash}p{(\linewidth - 4\tabcolsep) * \real{0.3333}}
  >{\raggedright\arraybackslash}p{(\linewidth - 4\tabcolsep) * \real{0.3333}}@{}}
\toprule\noalign{}
\begin{minipage}[b]{\linewidth}\raggedright
编号
\end{minipage} & \begin{minipage}[b]{\linewidth}\raggedright
公式
\end{minipage} & \begin{minipage}[b]{\linewidth}\raggedright
用途
\end{minipage} \\
\midrule\noalign{}
\endhead
\bottomrule\noalign{}
\endlastfoot
7-1 & \(F_s=\dfrac{\text{抗滑}}{\text{滑动}}\) & 定义 \\
7-2 & \(F_s=\dfrac{\tan\varphi}{\tan\beta}\) & 无黏性土坡 \\
7-3 & \(F_s=\dfrac{c_uLR}{Wd}\) & 整体圆弧滑动（φ=0 法） \\
7-4 & \(N_s=\dfrac{c}{\gamma H_{cr}}\) & 泰勒稳定数/极限坡高 \\
7-5 & 瑞典条分法
\(F_s=\dfrac{\sum\left[c_il_i+(W_i\cos\alpha_i-u_il_i)\tan\varphi_i\right]}{\sum W_i\sin\alpha_i}\)
& 一般圆弧滑动（总应力） \\
7-6 & 简化毕肖普
\(F_s=\dfrac{\sum\dfrac{c'b+(W-ub)\tan\varphi'}{\cos\alpha(1+\tan\alpha\tan\varphi'/F_s)}}{\sum W\sin\alpha}\)
& 有效应力条分法 \\
7-7 &
\(F_s=\dfrac{\gamma'}{\gamma_{sat}}\cdot\dfrac{\tan\varphi}{\tan\beta}\)
& 无黏性土坡顺坡渗流 \\
\end{longtable}
}

\begin{quote}
\textbf{考试高频提醒}：① 无黏性土坡
\(F_s=\tan\varphi/\tan\beta\)、只看坡角；② 条分法分母是
\(\sum W_i\sin\alpha_i\)，分子黏聚力项\textbf{乘以弧长}、摩擦项用有效法向力；③
简化毕肖普是隐式迭代，写清初值与收敛步骤；④ \(\alpha_i\)
符号别弄反------坡顶条为「下滑」，坡底条可抵消；⑤
有渗流必用有效应力指标与 \(u_i\)。
\end{quote}

\end{document}
